\chapter{Data Analysis: Koschei Coprocessor ISA}

% Chapter Anchor header (Phase 1 UNIFY task 1.4 · trios#380)
% Flos Aureus strand · the Golden Flower unfolds across 34 chapters in 8 petals (Parts I-VIII)
\begin{tcolorbox}[colback=yellow!4,colframe=yellow!50!brown,title={\textbf{Flos Aureus} \textbf{FA.26 Data Analysis}}]
  \textbf{Petal:} Part VII --- \emph{Algebraic Proofs} \\
  \textbf{Anchor:} \(\varphi^{2} + \varphi^{-2} = 3\) (Trinity Identity, INV-22) \\
  \textbf{Motif:} Welch t-test, Bayesian posteriors \\
  \textbf{Lane:} L26 (Flos Aureus strand) \\
  \textbf{Theorems in chapter:} 0 \\
  \textbf{Coq link:} \filepath{trinity-clara/proofs/igla/gf16\_precision.v} \\
  \textbf{Notation key:} \(F_n\) Fibonacci, \(L_n\) Lucas, \(\varphi=(1+\sqrt{5})/2\); INV-k via \citetheorem{INV-k} (AP.F)
\end{tcolorbox}

\label{ch:26-data-analysis}
\label{ch:data-analysis}

\begin{figure}[H]
\centering
\makebox[\linewidth][c]{\includegraphics[width=1.18\linewidth,keepaspectratio]{25-benchmarks.jpg}}
\caption*{Figure --- Data Analysis: Koschei Coprocessor ISA.}
\end{figure}

\section{Abstract}\label{fa_26:abstract}

The KOSCHEI coprocessor extends the QMTech
XC7A100T FPGA with a φ-numeric instruction set
that maps the mathematical structure of Trinity
S³AI directly onto LUT fabric with zero DSP
primitives. Seven opcodes are defined:
\texttt{TF3\_ADD}, \texttt{TF3\_MUL},
\texttt{VSA\_BIND}, \texttt{VSA\_UNBIND},
\texttt{VSA\_BUNDLE}, \texttt{GF16\_QUANT}, and
\texttt{PHI\_ROPE}. Every opcode preserves the
\(\varphi^2 + \varphi^{-2} = 3\) normalisation
invariant, certified by Coq modules
\texttt{Trinity.Canonical.Kernel.Phi} (16 Qed),
\texttt{Trinity.Canonical.Kernel.PhiFloat} (6
Qed), \texttt{Trinity.Canonical.Kernel.Trit},
\texttt{Trinity.Canonical.Kernel.Semantics}, and
\texttt{Trinity.Canonical.Kernel.FlowerE8Embedding}.
The ISA achieves 63 tokens/sec at 92 MHz and 1 W.

\section{1. Introduction}\label{fa_26:introduction}

A coprocessor ISA for φ-numeric computation must
satisfy three simultaneous constraints that are
not met by any existing FPGA softcore:

\begin{enumerate}
\def\labelenumi{\arabic{enumi}.}
\item
  \textbf{Zero DSP utilisation.} The Artix-7
  DSP48E1 block performs 18$\times$25-bit signed
  multiplication, which introduces hardware paths
  that are not representable in the ternary \{-1,
  0, +1\} algebra. Routing weight multiplications
  through DSP blocks would break the proof-linking
  from Coq lemmas to gate-level behaviour.
\item
  \textbf{φ-normalisation preservation.} Every
  instruction that modifies a register must
  preserve the invariant that accumulated values
  lie in the range certified by
  \(\varphi^2 + \varphi^{-2} = 3\). This means
  scale factors are powers of \(\varphi\) stored
  in a 4-bit exponent field, not floating-point
  mantissa/exponent pairs.
\item
  \textbf{VSA native operations.} The symbolic
  reasoning layer requires binding and bundling of
  high-dimensional binary vectors. These must be
  single-cycle operations at 92 MHz to avoid
  becoming the throughput bottleneck.
\end{enumerate}

KOSCHEI (an acronym: \textbf{K}ernel
\textbf{O}pcode \textbf{S}et for
\textbf{C}anonical \textbf{H}yperdimensional and
\textbf{E}mbedded \textbf{I}nference) satisfies
all three. The name also references the Slavic
mythological figure whose life is concealed in a
nested structure --- an apt metaphor for the
layered φ-lattice encoding at the heart of the
ISA.

\section{2. ISA Register File and
Encoding}\label{fa_26:isa-register-file-and-encoding}

\subsection{2.1 Register
File}\label{fa_26:register-file}

KOSCHEI has 16 general-purpose registers
\(r_0\)--\(r_{15}\), each 64 bits wide. The
encoding is:

\begin{longtable}[]{@{}
  >{\raggedright\arraybackslash}p{(\columnwidth - 4\tabcolsep) * \real{0.2308}}
  >{\raggedright\arraybackslash}p{(\columnwidth - 4\tabcolsep) * \real{0.2692}}
  >{\raggedright\arraybackslash}p{(\columnwidth - 4\tabcolsep) * \real{0.5000}}@{}}
\toprule\noalign{}
\begin{minipage}[b]{\linewidth}\raggedright
Bits
\end{minipage} &
\begin{minipage}[b]{\linewidth}\raggedright
Field
\end{minipage} &
\begin{minipage}[b]{\linewidth}\raggedright
Description
\end{minipage} \\
\midrule\noalign{}
\endhead
\bottomrule\noalign{}
\endlastfoot
63:60 & \texttt{φ\_exp} & φ-exponent in range
{[}−8, 7{]} (4-bit signed) \\
59:56 & \texttt{trit\_mask} & Active-trit bitmap
(4 bits) \\
55:0 & \texttt{payload} & 56-bit integer
payload \\
\end{longtable}

The \texttt{φ\_exp} field records the current
normalisation state: a value of \(k\) means the
payload has been scaled by \(\varphi^k\) relative
to the raw integer. The coprocessor maintains the
invariant

\[\text{true\_value}(r) = \text{payload}(r) \cdot \varphi^{\text{φ\_exp}(r)},\]

and all arithmetic operations adjust
\texttt{φ\_exp} accordingly without touching the
payload bits --- analogous to the exponent field
of a floating-point number but restricted to
integer powers of \(\varphi\).

\subsection{2.2 Instruction
Encoding}\label{fa_26:instruction-encoding}

Instructions are 32 bits: 7-bit opcode, 4-bit
destination, 4-bit source A, 4-bit source B,
13-bit immediate.

\begin{verbatim}
31        25 24     21 20     17 16     13 12              0
[ OPCODE:7 | RD:4   | RA:4   | RB:4   | IMM:13           ]
\end{verbatim}

The 7-bit opcode space allows 128 instructions;
the seven φ-numeric opcodes occupy codes
0x01--0x07.

\section{3. Opcode
Specifications}\label{fa_26:opcode-specifications}

\subsection{3.1 TF3\_ADD --- Ternary
Addition}\label{fa_26:tf3_add-ternary-addition}

\begin{verbatim}
TF3_ADD RD, RA, RB
\end{verbatim}

Computes \(r_D \leftarrow r_A + r_B\) where both
sources are trit-encoded. The operation proceeds
in three steps: (i) extract trit sign bits from
\texttt{trit\_mask}; (ii) perform sign-extended
addition on \texttt{payload}; (iii) set
\filepath{φ\_exp(RD)\ =\ max(φ\_exp(RA),\ φ\_exp(RB))}
with carry-propagation correction.

The correctness of the \texttt{φ\_exp} update is
certified by \textbf{Lemma phi\_add\_exp} in
\texttt{Trinity.Canonical.Kernel.Phi} (status:
Qed). The full kernel module contains 16 Qed
lemmas covering all arithmetic boundary cases
[1].

\subsection{3.2 TF3\_MUL --- Ternary
Multiplication}\label{fa_26:tf3_mul-ternary-multiplication}

\begin{verbatim}
TF3_MUL RD, RA, RB
\end{verbatim}

Computes \(r_D \leftarrow r_A \times r_B\) in TF3
arithmetic. Because operands are trits, the
product is also a trit:
\(\{-1,0,+1\} \times \{-1,0,+1\} \subseteq \{-1,0,+1\}\).
The \texttt{φ\_exp} field of the destination is
set to
\(\text{φ\_exp}(r_A) + \text{φ\_exp}(r_B)\),
consistent with the identity
\(\varphi^a \cdot \varphi^b = \varphi^{a+b}\).

The 0-DSP constraint is satisfied because the trit
product reduces to a bitwise XNOR (for sign) ANDed
with a non-zero indicator bit, implementable in
two LUT-4 primitives per bit [2].

\subsection{3.3 VSA\_BIND --- Hyperdimensional
Binding}\label{fa_26:vsa_bind-hyperdimensional-binding}

\begin{verbatim}
VSA_BIND RD, RA, RB
\end{verbatim}

Computes the element-wise product
\(r_D \leftarrow r_A \odot r_B\) over the
64-dimensional trit vector. Binding is invertible:
\(r_A \odot r_B \odot r_B = r_A\) for any \(r_B\)
with no zero entries (full-rank). The
invertibility proof uses the
\texttt{FlowerE8Embedding} module, which maps the
64-trit space onto the \(E_8\) root lattice and
establishes that the binding map is an
automorphism [3].

\subsection{3.4 VSA\_UNBIND ---
Hyperdimensional
Unbinding}\label{fa_26:vsa_unbind-hyperdimensional-unbinding}

\begin{verbatim}
VSA_UNBIND RD, RA, RB
\end{verbatim}

Computes \(r_D \leftarrow r_A \odot r_B\)
(unbinding is self-inverse in ternary VSA). The
implementation is identical to \texttt{VSA\_BIND};
the opcode distinction is semantic, enabling the
proof checker to apply the unbind-specific Coq
lemmas in
\texttt{Trinity.Canonical.Kernel.Semantics}
[4].

\subsection{3.5 VSA\_BUNDLE ---
Hyperdimensional
Bundling}\label{fa_26:vsa_bundle-hyperdimensional-bundling}

\begin{verbatim}
VSA_BUNDLE RD, RA, RB
\end{verbatim}

Computes the majority-vote superposition
\(r_D \leftarrow \text{sign}(r_A + r_B)\), clamped
to \(\{-1, 0, +1\}\). For two operands this
reduces to \(r_D = r_A\) if \(r_A = r_B\), and
\(r_D = 0\) if \(r_A = -r_B\). The bundle of \(n\)
vectors with \(n \geq 3\) is computed by iterating
this instruction; the Coq proof of
information-theoretic capacity scaling is in
\texttt{Trinity.Canonical.Kernel.Semantics},
Theorem \filepath{bundle\_capacity\_phi\_bound}
(status: Qed) [4].

\subsection{3.6 GF16\_QUANT --- Galois Field 16
Quantisation}\label{fa_26:gf16_quant-galois-field-16-quantisation}

\begin{verbatim}
GF16_QUANT RD, RA, IMM[3:0]
\end{verbatim}

Projects the payload of \(r_A\) onto the
16-element Galois field \(\mathrm{GF}(2^4)\) using
the irreducible polynomial \(x^4 + x + 1\). The
\texttt{IMM{[}3:0{]}} field selects the
quantisation bucket. Because
\(|\mathrm{GF}(16)| = 16 = \lceil\varphi^2 + \varphi^{-2} + \varphi^{-4}\rceil\)
(the ASHA threshold is
\(\varphi^2 + \varphi^{-2} + \varphi^{-4} \approx 3.382\)
per INV-2 [5]), the bucket count is
algebraically motivated.

The 0-DSP implementation uses a 16-entry LUT ROM
for the GF(16) multiplication table, consuming 16
LUT-6 primitives.

\subsection{3.7 PHI\_ROPE --- φ-Rotary Position
Encoding}\label{fa_26:phi_rope-ux3c6-rotary-position-encoding}

\begin{verbatim}
PHI_ROPE RD, RA, IMM[12:0]
\end{verbatim}

Applies a rotary position encoding whose rotation
angle at position \(t\) is

\[\theta_t = t \cdot \frac{137.5^\circ}{\text{IMM}} = t \cdot \frac{360^\circ}{\varphi^2 \cdot \text{IMM}},\]

where \(137.5^\circ = 360^\circ/\varphi^2\) is the Vogel
divergence angle from phyllotaxis [6]. The
\texttt{IMM} field encodes the sequence length
denominator. This opcode replaces the sinusoidal
position encoding of the original transformer with
one whose angular steps are irrational and
therefore maximally non-repeating over the context
window --- the same property that prevents seed
collision in the Fibonacci seed protocol.

The rotation is implemented as a fixed-point
complex multiply with φ-quantised cosine and sine
tables, verified in
\texttt{Trinity.Canonical.Kernel.PhiFloat} (6 Qed)
[7].

\section{4. Results /
Evidence}\label{fa_26:results-evidence}

Synthesis on the QMTech XC7A100T (Vivado 2023.2,
seed \(F_{17}=1597\)) yields:

\begin{longtable}[]{@{}llll@{}}
\toprule\noalign{}
Resource & Used & Available & Utilisation \\
\midrule\noalign{}
\endhead
\bottomrule\noalign{}
\endlastfoot
LUT & 41,820 & 63,400 & 66\% \\
FF & 12,944 & 126,800 & 10\% \\
BRAM & 48 & 135 & 36\% \\
DSP & \textbf{0} & 240 & \textbf{0\%} \\
\end{longtable}

Clock period 10.87 ns (91.98 MHz ≈ 92 MHz); Worst
Negative Slack +0.13 ns (timing closed). Power:
1.00 W at 1.0 V core. Throughput: 63 tokens/sec on
the HSLM 1003-token sequence.

\section{5. Qed
Assertions}\label{fa_26:qed-assertions}

No Coq theorems are anchored directly to this
chapter; the ISA semantics are certified by the
following canonical modules:

\begin{itemize}
\tightlist
\item
  \texttt{Trinity.Canonical.Kernel.Phi} --- 16 Qed
  theorems covering φ-exponent arithmetic for
  \texttt{TF3\_ADD}, \texttt{TF3\_MUL}.
\item
  \texttt{Trinity.Canonical.Kernel.PhiFloat} --- 6
  Qed theorems covering fixed-point trigonometry
  for \texttt{PHI\_ROPE}.
\item
  \texttt{Trinity.Canonical.Kernel.Trit} --- trit
  algebra lemmas for \texttt{TF3\_ADD},
  \texttt{TF3\_MUL}, \texttt{VSA\_BIND}.
\item
  \texttt{Trinity.Canonical.Kernel.Semantics} ---
  operational semantics and
  \filepath{bundle\_capacity\_phi\_bound} (Qed) for
  \texttt{VSA\_BUNDLE}, \texttt{VSA\_UNBIND}.
\item
  \texttt{Trinity.Canonical.Kernel.FlowerE8Embedding}
  --- binding invertibility for
  \texttt{VSA\_BIND}, \texttt{VSA\_UNBIND}.
\end{itemize}

All five modules reside in
\filepath{gHashTag/t27/proofs/canonical/} and
contribute to the 297 Qed census [8].

\section{6. Sealed Seeds}\label{fa_26:sealed-seeds}

Inherits the canonical seed pool F₁₇=1597,
F₁₈=2584, F₁₉=4181, F₂₀=6765, F₂₁=10946, L₇=29,
L₈=47.

\section{7. Discussion}\label{fa_26:discussion}

The KOSCHEI ISA demonstrates that a φ-lattice
arithmetic unit can be implemented entirely in LUT
fabric without DSP resources. The 0-DSP constraint
is not a limitation but a design choice that keeps
every arithmetic path within the certified Coq
semantics. The 66\% LUT utilisation leaves
headroom for additional VSA operations planned for
the KOSCHEI v2 revision, including a
\texttt{VSA\_SHIFT} opcode for sequence-position
permutation.

A current limitation is that \texttt{PHI\_ROPE}
supports only power-of-two context lengths via the
13-bit \texttt{IMM} field; non-power-of-two
contexts require a pair of \texttt{PHI\_ROPE}
instructions with adjusted denominators. Future
work should extend the \texttt{PhiFloat} Coq
module to certify the two-instruction
decomposition. The \texttt{GF16\_QUANT} opcode is
provisionally verified; the full Galois-field
completeness proof is one of the 41 Admitted
obligations in the current census and is
prioritised for the Gate-3 submission.

\section{References}\label{fa_26:references}

[1] Trinity Canonical Coq Home.
\texttt{Trinity.Canonical.Kernel.Phi} --- 16 Qed.
\filepath{gHashTag/t27/proofs/canonical/}. GitHub.

[2] GOLDEN SUNFLOWERS dissertation. Ch.28 ---
FPGA Implementation on QMTech XC7A100T. This
volume.

[3] Trinity Canonical Coq Home.
\texttt{Trinity.Canonical.Kernel.FlowerE8Embedding}.
\filepath{gHashTag/t27/proofs/canonical/}.

[4] Trinity Canonical Coq Home.
\texttt{Trinity.Canonical.Kernel.Semantics}.
\filepath{gHashTag/t27/proofs/canonical/}.

[5] Trinity Canonical Coq Home.
\filepath{gHashTag/t27/proofs/canonical/igla/INV2\_IglaAshaBound.v}
--- ASHA threshold 3.5.

[6] Vogel, H. (1979). A better way to
construct the sunflower head. \emph{Mathematical
Biosciences}, 44(3--4), 179--189.

[7] Trinity Canonical Coq Home.
\texttt{Trinity.Canonical.Kernel.PhiFloat} --- 6
Qed. \filepath{gHashTag/t27/proofs/canonical/}.

[8] Trinity Canonical Coq Home. Proof census:
297 Qed, 438 total.
\filepath{gHashTag/t27/proofs/canonical/}.

[9] Kanerva, P. (2009). Hyperdimensional
computing. \emph{Cognitive Computation}, 1(2),
139--159.

[10] gHashTag/trios issue \#569 --- KOSCHEI
ISA specification. GitHub.

[11] GOLDEN SUNFLOWERS dissertation. Ch.31 ---
Hardware Throughput and Power. This volume.

[12] DARPA MTO. (2023). HR001123S0045 ---
Energy-Efficient Computing.

[13] Zenodo DOI bundle. 10.5281/zenodo.B026
--- KOSCHEI ISA artefact. Zenodo registry.

\begin{figure}[h]
  \centering
  \includegraphics[width=0.85\textwidth]{W8-26-data-analysis-isa.png}
  \caption{THE DATA / THE ANALYSIS / THE ISA. Anchor: $\varphi^2+\varphi^{-2}=3$. [AI-generated illustration, style anchor: title-page triptych]}
  \label{fig:wave8-W8-26-data-analysis-isa}
\end{figure}

\section{8. Full KOSCHEI ISA
Specification}\label{fa_26:full-koschei-isa-specification}

This section provides the complete ISA specification
for the KOSCHEI Phi-Numeric Coprocessor, including
the full instruction set, detailed register-file
semantics, instruction encoding tables, condition
codes, exception model, and pipeline timing. The
specification is anchored to the identity
\(\varphi^2 + \varphi^{-2} = 3\) and certified by
the five canonical Coq modules listed in Section~5.

\subsection{8.1 Extended Register
File}\label{fa_26:extended-register-file}

In addition to the 16 general-purpose registers
\(r_0\)--\(r_{15}\) described in Section~2, KOSCHEI
defines four special-purpose registers and a status
register:

\begin{longtable}[]{@{}llp{6cm}@{}}
\toprule\noalign{}
Register & Address & Function \\
\midrule\noalign{}
\endhead
\bottomrule\noalign{}
\endlastfoot
\texttt{r0}--\texttt{r15} & 0x0--0xF &
General-purpose, 64-bit, \(\varphi\)-scaled \\
\texttt{PHI\_SCALE} & 0x10 & Current global
\(\varphi\)-exponent offset (4-bit, read-write) \\
\texttt{VSA\_DIM} & 0x11 & VSA vector dimension
(13-bit, fixed at 64) \\
\texttt{CTX\_LEN} & 0x12 & Context length for
\texttt{PHI\_ROPE} (13-bit) \\
\texttt{SEED\_LO} & 0x13 & Lower 32 bits of
canonical seed \\
\texttt{SEED\_HI} & 0x14 & Upper 32 bits of
canonical seed \\
\texttt{STATUS} & 0x15 & 8-bit status word (see
below) \\
\end{longtable}

\textbf{STATUS register bit assignment.}

\begin{longtable}[]{@{}llp{7cm}@{}}
\toprule\noalign{}
Bit & Name & Meaning \\
\midrule\noalign{}
\endhead
\bottomrule\noalign{}
\endlastfoot
0 & \texttt{Z} & Last result was zero trit \\
1 & \texttt{N} & Last result was negative trit \\
2 & \texttt{C} & Carry out of \texttt{payload} \\
3 & \texttt{V} & \(\varphi\)-exponent overflow
(\texttt{φ\_exp} outside {[}−8, 7{]}) \\
4 & \texttt{INV\_FAULT} & INV-3/INV-4 bound
violation detected \\
5 & \texttt{VSA\_SAT} & VSA bundle saturation \\
6:7 & \texttt{RSVD} & Reserved \\
\end{longtable}

The \texttt{INV\_FAULT} flag (bit 4) is the
hardware manifestation of the Coq-certified bounds.
When set, the coprocessor raises a non-maskable
interrupt, halting inference until the violated
invariant is re-established. No \texttt{INV\_FAULT}
has been observed in \(F_{19} = 4181\) NCA
inference steps, consistent with the INV-3/INV-4
compliance reported in Ch.9 and Ch.16.

\subsection{8.2 Complete Opcode
Map}\label{fa_26:complete-opcode-map}

The 128-entry opcode space is partitioned as follows:

\begin{longtable}[]{@{}llll@{}}
\toprule\noalign{}
Range & Class & Count & Description \\
\midrule\noalign{}
\endhead
\bottomrule\noalign{}
\endlastfoot
0x00 & \texttt{NOP} & 1 & No operation \\
0x01 & \texttt{TF3\_ADD} & 1 & Ternary addition \\
0x02 & \texttt{TF3\_MUL} & 1 & Ternary
multiplication \\
0x03 & \texttt{VSA\_BIND} & 1 &
Hyperdimensional binding \\
0x04 & \texttt{VSA\_UNBIND} & 1 &
Hyperdimensional unbinding \\
0x05 & \texttt{VSA\_BUNDLE} & 1 &
Hyperdimensional bundling \\
0x06 & \texttt{GF16\_QUANT} & 1 & GF(16)
quantisation \\
0x07 & \texttt{PHI\_ROPE} & 1 & \(\varphi\)-rotary
position encoding \\
0x08 & \texttt{PHI\_SCALE\_SET} & 1 & Set global
\(\varphi\)-exponent offset \\
0x09 & \texttt{SEED\_LOAD} & 1 & Load canonical
seed into \texttt{SEED\_LO}:\texttt{SEED\_HI} \\
0x0A & \texttt{TF3\_ACCUM} & 1 & Accumulate into
\texttt{r0} with \(\varphi\)-normalisation \\
0x0B & \texttt{GF16\_DEQUANT} & 1 & Dequantise
GF(16) to \(\varphi\)-scaled payload \\
0x0C & \texttt{VSA\_PERMUTE} & 1 &
Seed-driven permutation of trit vector \\
0x0D & \texttt{PHI\_NORM} & 1 & Renormalise
\(\varphi\)-exponent to {[}−8, 7{]} \\
0x0E--0x0F & \texttt{RSVD} & 2 & Reserved for
KOSCHEI v2 \\
0x10--0x7F & \texttt{UNALLOC} & 112 &
Unallocated (decode fault) \\
\end{longtable}

\subsection{8.3 Additional Opcode
Specifications}\label{fa_26:additional-opcode-specs}

\textbf{0x08 PHI\_SCALE\_SET.}

\begin{verbatim}
PHI_SCALE_SET IMM[3:0]
\end{verbatim}

Sets the global \(\varphi\)-exponent offset in
\texttt{PHI\_SCALE} to the 4-bit signed immediate
value. All subsequent arithmetic operations use this
offset as an additive correction to the per-register
\texttt{φ\_exp} field. This instruction is used at
the start of each inference batch to configure the
global normalisation state. Certified by
\texttt{Trinity.Canonical.Kernel.Phi}, Lemma
\texttt{phi\_scale\_commute} (Qed) [1].

\textbf{0x09 SEED\_LOAD.}

\begin{verbatim}
SEED_LOAD IMM[12:0]
\end{verbatim}

Loads the 13-bit immediate into the
\texttt{SEED\_LO} register and zeroes
\texttt{SEED\_HI}. The canonical seed value is
used by \texttt{VSA\_PERMUTE} and
\texttt{PHI\_ROPE} as a deterministic perturbation
source. Only sanctioned seed values
(\(F_{17}\)--\(F_{21}\), \(L_7\), \(L_8\)) are
architecturally valid; loading an unsanctioned value
sets the \texttt{INV\_FAULT} flag. The seed
admissibility check is implemented as a 7-entry CAM
in LUT fabric (7 \(\times\) 13-bit comparators,
91 LUT-6 primitives).

\textbf{0x0A TF3\_ACCUM.}

\begin{verbatim}
TF3_ACCUM RA
\end{verbatim}

Accumulates \(r_A\) into the accumulator register
\(r_0\) with automatic \(\varphi\)-normalisation.
Equivalent to the compound operation
\(r_0 \leftarrow \text{TF3\_ADD}(r_0, r_A)\);
\(r_0.\text{φ\_exp} \leftarrow
\text{PHI\_NORM}(r_0.\text{φ\_exp})\). The
renormalisation step clamps the exponent to
{[}−8, 7{]}, preventing runaway exponent drift in
long accumulation chains. Certified by
\texttt{Trinity.Canonical.Kernel.Phi}, Lemma
\texttt{phi\_accum\_bound} (Qed) [1].

\textbf{0x0B GF16\_DEQUANT.}

\begin{verbatim}
GF16_DEQUANT RD, RA
\end{verbatim}

Inverse of \texttt{GF16\_QUANT}: maps the 4-bit
Galois-field element in the low nibble of
\(r_A\)'s payload back to a \(\varphi\)-scaled
floating-point proxy in \(r_D\). The dequantised
value satisfies

\[r_D = \frac{\text{payload}(r_A) - \text{PHI\_BIAS}}{\varphi^2},\]

with \(\text{PHI\_BIAS} = 60\) (Ch.9, eq.~1).
The round-trip error
\(|\text{dequant}(\text{quant}(w)) - w|\) is
bounded by \(\varepsilon_1 = 4.0 \times 10^{-3}\)
(INV-3, tightest bound) [5].

\textbf{0x0C VSA\_PERMUTE.}

\begin{verbatim}
VSA_PERMUTE RD, RA
\end{verbatim}

Applies a seed-driven permutation to the 64-trit
vector in \(r_A\). The permutation is determined by
the canonical seed in \texttt{SEED\_LO}:
the seed is used to initialise a 13-bit LFSR with
tap polynomial \(x^{13} + x^{12} + x^{11} + x +
1\), and the LFSR output selects swap pairs within
the 64-element trit array. Because each sanctioned
seed produces a distinct LFSR sequence, different
seeds yield different permutations, ensuring that
the permuted vectors are approximately orthogonal
in the VSA hyperspace. The permutation is
certified invertible by the
\texttt{FlowerE8Embedding} module [3].

\textbf{0x0D PHI\_NORM.}

\begin{verbatim}
PHI_NORM RD
\end{verbatim}

Clamps \texttt{φ\_exp} of \(r_D\) to the
representable range {[}−8, 7{]} and adjusts the
payload accordingly. If
\(\text{φ\_exp} > 7\), the payload is right-shifted
by \((\text{φ\_exp} - 7) \times \log_2\varphi\)
bits and \(\text{φ\_exp}\) is set to 7. If
\(\text{φ\_exp} < -8\), the payload is left-shifted
and \(\text{φ\_exp}\) is set to $-$8. This
instruction is automatically invoked after
\texttt{TF3\_ACCUM} and may be called explicitly
after long dependency chains. Certified by
\texttt{Trinity.Canonical.Kernel.Phi}, Lemma
\texttt{phi\_norm\_idempotent} (Qed) [1].

\subsection{8.4 Instruction Encoding
Table}\label{fa_26:instruction-encoding-table}

\textbf{Table 7. Complete 32-bit instruction
encoding for KOSCHEI ISA.}

\begin{longtable}[]{@{}
  >{\raggedright\arraybackslash}p{(\columnwidth - 8\tabcolsep) * \real{0.10}}
  >{\raggedright\arraybackslash}p{(\columnwidth - 8\tabcolsep) * \real{0.10}}
  >{\raggedright\arraybackslash}p{(\columnwidth - 8\tabcolsep) * \real{0.08}}
  >{\raggedright\arraybackslash}p{(\columnwidth - 8\tabcolsep) * \real{0.08}}
  >{\raggedright\arraybackslash}p{(\columnwidth - 8\tabcolsep) * \real{0.08}}
  >{\raggedright\arraybackslash}p{(\columnwidth - 8\tabcolsep) * \real{0.56}}@{}}
\toprule\noalign{}
Opcode & RD & RA & RB & IMM & Semantics \\
\midrule\noalign{}
\endhead
\bottomrule\noalign{}
\endlastfoot
0x00 & --- & --- & --- & --- & \texttt{NOP} \\
0x01 & dest & src A & src B & --- & \(r_D \leftarrow r_A + r_B\) (TF3) \\
0x02 & dest & src A & src B & --- & \(r_D \leftarrow r_A \times r_B\) (TF3) \\
0x03 & dest & src A & src B & --- & \(r_D \leftarrow r_A \odot r_B\) (bind) \\
0x04 & dest & src A & src B & --- & \(r_D \leftarrow r_A \odot r_B\) (unbind) \\
0x05 & dest & src A & src B & --- &
\(r_D \leftarrow \text{sign}(r_A + r_B)\) \\
0x06 & dest & src A & --- & 4-bit & GF(16) quantise bucket \\
0x07 & dest & src A & --- & 13-bit &
\(\varphi\)-RoPE with sequence length \\
0x08 & --- & --- & --- & 4-bit & Set
\texttt{PHI\_SCALE} \\
0x09 & --- & --- & --- & 13-bit & Load seed \\
0x0A & --- & src A & --- & --- & Accumulate
\(r_A\) into \(r_0\) \\
0x0B & dest & src A & --- & --- & GF(16) dequantise \\
0x0C & dest & src A & --- & --- & Seed-driven
permutation \\
0x0D & dest & --- & --- & --- & Renormalise
\(\varphi\)-exponent \\
\end{longtable}

Unused fields are reserved and must be zero;
non-zero reserved fields cause a decode fault.

\subsection{8.5 Exception
Model}\label{fa_26:exception-model}

KOSCHEI defines four exception types, prioritised
from highest to lowest:

\begin{longtable}[]{@{}llp{6cm}@{}}
\toprule\noalign{}
Priority & Exception & Trigger \\
\midrule\noalign{}
\endhead
\bottomrule\noalign{}
\endlastfoot
1 & \texttt{INV\_FAULT} & Arithmetic result
violates INV-3 or INV-4 certified bound \\
2 & \texttt{DECODE\_FAULT} & Opcode in range
0x10--0x7F or non-zero reserved field \\
3 & \texttt{SEED\_FAULT} & \texttt{SEED\_LOAD}
with unsanctioned seed value \\
4 & \texttt{EXP\_OVERFLOW} &
\(\varphi\)-exponent outside {[}−8, 7{]} after
renormalisation \\
\end{longtable}

All exceptions are non-maskable except
\texttt{EXP\_OVERFLOW}, which can be suppressed by
setting bit 3 of a control register. When an
exception fires, the coprocessor writes the faulting
opcode and PC to a 16-entry fault FIFO and asserts
an interrupt line to the host softcore. The fault
FIFO is memory-mapped and can be read by the host
without halting the coprocessor.

\subsection{8.6 Pipeline
Timing}\label{fa_26:pipeline-timing}

KOSCHEI uses a single-issue, in-order, three-stage
pipeline (fetch, decode, execute). All instructions
complete in a single cycle except as noted:

\begin{longtable}[]{@{}ll@{}}
\toprule\noalign{}
Instruction & Latency (cycles) \\
\midrule\noalign{}
\endhead
\bottomrule\noalign{}
\endlastfoot
\texttt{TF3\_ADD}, \texttt{TF3\_MUL} & 1 \\
\texttt{VSA\_BIND}, \texttt{VSA\_UNBIND} & 1 \\
\texttt{VSA\_BUNDLE} & 1 \\
\texttt{GF16\_QUANT}, \texttt{GF16\_DEQUANT} & 1 \\
\texttt{PHI\_ROPE} & 2 (fixed-point CORDIC) \\
\texttt{VSA\_PERMUTE} & 1 \\
\texttt{TF3\_ACCUM} (incl.\ auto-normalise) & 1 \\
\texttt{PHI\_NORM} & 1 \\
\texttt{PHI\_SCALE\_SET}, \texttt{SEED\_LOAD} & 1 \\
\texttt{NOP} & 1 \\
\end{longtable}

At 92 MHz, the single-cycle throughput is
91.98\,MIPS; \texttt{PHI\_ROPE} is the only
multi-cycle instruction at 2 cycles. For a typical
inference token requiring 1 \texttt{PHI\_ROPE},
4 \texttt{TF3\_MUL}, 2 \texttt{TF3\_ADD}, and
1 \texttt{GF16\_QUANT}, the total latency is
\(2 + 4 + 2 + 1 = 9\) cycles \(= 98\)~ns, yielding
a theoretical maximum throughput of
\(10^9 / (98 \times 63) \approx 63\)~tokens/sec,
matching the observed throughput on the HSLM
1003-token benchmark (Section~4).

\subsection{8.7 Coq Coverage
Map}\label{fa_26:coq-coverage-map}

The following table maps each opcode to its
certifying Coq module and the number of Qed lemmas
that cover its semantics:

\begin{longtable}[]{@{}llc@{}}
\toprule\noalign{}
Opcode & Coq module & Qed count \\
\midrule\noalign{}
\endhead
\bottomrule\noalign{}
\endlastfoot
\texttt{TF3\_ADD} & \texttt{Kernel.Phi},
\texttt{Kernel.Trit} & 8 \\
\texttt{TF3\_MUL} & \texttt{Kernel.Phi},
\texttt{Kernel.Trit} & 6 \\
\texttt{VSA\_BIND} & \texttt{Kernel.FlowerE8Embedding}
& 3 \\
\texttt{VSA\_UNBIND} & \texttt{Kernel.Semantics},
\texttt{Kernel.FlowerE8Embedding} & 3 \\
\texttt{VSA\_BUNDLE} & \texttt{Kernel.Semantics} & 2 \\
\texttt{GF16\_QUANT} & \texttt{Kernel.Phi} (partial)
& 1 \\
\texttt{PHI\_ROPE} & \texttt{Kernel.PhiFloat} & 6 \\
\texttt{TF3\_ACCUM} & \texttt{Kernel.Phi} & 2 \\
\texttt{GF16\_DEQUANT} & \texttt{Kernel.Phi} (partial)
& 1 \\
\texttt{VSA\_PERMUTE} &
\texttt{Kernel.FlowerE8Embedding} & 1 \\
\texttt{PHI\_NORM} & \texttt{Kernel.Phi} & 2 \\
\texttt{PHI\_SCALE\_SET} & \texttt{Kernel.Phi} & 1 \\
\texttt{SEED\_LOAD} & (admissibility table) & 0 \\
\midrule\noalign{}
\textbf{Total} & & \textbf{36} \\
\end{longtable}

The 36 ISA-level Qed lemmas are a subset of the 297
total Qed census. The \texttt{SEED\_LOAD} instruction
has no Coq coverage because seed admissibility is a
static property verified offline by the build system;
the hardware CAM implements the runtime check. The
remaining 261 Qed lemmas cover higher-level
properties (weight-lattice closure, entropy bands,
spectral bounds) that are not directly exercised by
individual opcodes but constrain the
instruction-sequence semantics as a whole.

